% mazzi: 09,10
\chapter{L'apparato scheletrico}

% ---------------------------------------------------------------------------
\section{Funzioni dell'apparato scheletrico}

\textbf{Apparato scheletrico}: \textbf{protezione}, \textbf{forma} e \textbf{sostegno},
\textbf{movimento}, \textbf{riserva di minerali}, \textbf{emopoiesi}. Si distingue in
\textbf{appendicolare} e \textbf{assile}.
\begin{itemize}
  \item $\sim$20\% del peso corporeo; nell'adulto $\sim$206 ossa (esclusi ossicini dell'udito e
    sesamoidi);
  \item conformazione modellata dalla \textbf{forza di gravità} e dall'uso di arti inferiori e
    superiori.
\end{itemize}

% ---------------------------------------------------------------------------
\section{Evoluzione della postura eretta}

Passaggio da \textbf{quadrupede} a \textbf{bipede} $\rightarrow$ rimodellamento del \textbf{bacino}:
\begin{itemize}
  \item mammiferi inferiori: bacino stretto sul piano frontale, allungato in senso sagittale;
  \item uomo: asse laterolaterale maggiore, asse sagittale minore.
\end{itemize}
Evoluzione del cranio: scimpanzé $\rightarrow$ \textit{Australopithecus africanus} (3,7 mln di anni)
$\rightarrow$ \textit{Homo abilis} (2,2 mln) $\rightarrow$ \textit{Homo erectus} (1 mln) $\rightarrow$
\textit{Homo sapiens} (da 200.000 anni) $\rightarrow$ aumento della capacità cranica e riduzione del
massiccio facciale.

% ---------------------------------------------------------------------------
\section{Equilibrio corporeo e arti}

\begin{itemize}
  \item \textbf{arti inferiori}: deambulazione ed equilibrio statico/dinamico; più robusti dei
    superiori;
  \item \textbf{arti superiori}: equilibrio, deambulazione e varietà di azioni;
  \item la \textbf{colonna vertebrale} presenta curvature opposte (di lato) $\rightarrow$
    riequilibrio del peso e ammortizzazione della gravità su cranio e tronco.
\end{itemize}

\rubrica{Veduta anteriore dello scheletro}
\begin{itemize}
  \item \textbf{cranio}: orbita, mascella, mandibola, articolazione temporo-mandibolare;
  \item \textbf{cingolo scapolare}: clavicola (articolazione sternoclavicolare) e scapola (processo
    coracoideo, acromion);
  \item \textbf{arto superiore}: omero (piccola e grande tuberosità); avambraccio = radio + ulna
    (articolazioni radioulnari prossimale e distale); mano = carpo, metacarpo, dita;
  \item \textbf{sterno}: manubrio, corpo, processo xifoideo;
  \item \textbf{bacino}: iliaco, pubico, ischiatico (sinfisi pubica; sacro e coccige); anca
    (coxo-femorale);
  \item \textbf{arto inferiore}: femore; ginocchio (femoro-patellare, patella); gamba = tibia
    (tuberosità tibiale) + perone; piede = astragalo, ossa tarsali, navicolare, cuneiformi, cuboide,
    metatarso, falangi.
\end{itemize}

\rubrica{Veduta posteriore}
\begin{itemize}
  \item cranio: parietale e occipitale; atlante ed epistrofeo;
  \item spalla: spina della scapola, articolazione acromioclavicolare, testa dell'omero nella
    cavità glenoidea;
  \item colonna vertebrale; gomito (olecrano dell'ulna); cresta iliaca; coxofemorale (collo del
    femore, acetabolo);
  \item carpo: scafoide, semilunare, piramidale, uncinato, capitato, trapezoide, trapezio;
  \item femore: testa, piccolo e grande trocantere; ginocchio: condilo mediale e laterale;
  \item gamba/caviglia: articolazione tibioperoneale, testa del perone, piatto tibiale, malleolo
    mediale, sindesmosi tibioperoneale, malleolo laterale, articolazione talocrurale, astragalo,
    subtalare, calcagno.
\end{itemize}

% ---------------------------------------------------------------------------
\section{Il cranio}

\begin{itemize}
  \item articolazioni tra le ossa craniche/facciali = tutte \textbf{sinartrosi} (per lo più
    \textbf{suture}); una sola mobile = \textbf{diartrosi} tra mandibola e osso temporale;
  \item \textbf{neurocranio}: protegge l'encefalo;
  \item \textbf{splancnocranio} (massiccio facciale): in rapporto con la porzione iniziale di
    apparato respiratorio e digerente.
\end{itemize}

\subsection{Neurocranio e splancnocranio}
\begin{itemize}
  \item \textbf{neurocranio}: \textbf{base} e \textbf{volta}, ossa piatte e brevi $\rightarrow$
    occipitale, sfenoide, temporale, etmoide, frontale, parietale;
  \item \textbf{splancnocranio}: lacrimale, nasale, zigomatico, mascellare, mandibola.
\end{itemize}
Suture principali: sagittale (tra i parietali), lambdoidea (parietali--occipitale), coronale
(parietali--frontale), squamosa (parietale--temporale). Elementi della base cranica (visione
inferiore, mandibola rimossa): forame palatino maggiore, fossa incisiva, palatino, vomere, processi
pterigoidei dello sfenoide, arco zigomatico, forame ovale, forame lacero, canale carotideo, meato
acustico esterno, fossa mandibolare, forame giugulare, forame stilomastoideo, forame magno, linea
nucale superiore.

\subsection{Osso occipitale}
Impari, a conchiglia/romboidale, base del neurocranio.
\begin{itemize}
  \item \textbf{forame magno} ($\phi$ medio $\sim$35 mm): passaggio di midollo spinale, arterie
    vertebrali, radici spinali, XI paio dei nervi cranici; il bulbo vi si estende nel canale
    vertebrale $\rightarrow$ midollo spinale;
  \item si articola anteriormente per sincondrosi (via clivo) col corpo dello sfenoide.
\end{itemize}

\subsection{Osso sfenoide}
Impari e mediano, al centro della base cranica.
\begin{itemize}
  \item articolazioni: etmoide e frontale (avanti), occipite (dietro), temporale e parietale (lati);
  \item partecipa a orbite, fosse nasali, fossa temporale e sotto-temporale;
  \item corpo cuboidale a 6 facce (2 laterali + superiore endocraniche; anteriore e inferiore
    esocraniche; posteriore per sinostosi col clivo); contiene i \textbf{seni sfenoidali}
    (comunicanti con le cavità nasali).
\end{itemize}

\subsection{Osso temporale}
Pari; anteriore all'occipitale, posteriore alle grandi ali dello sfenoide, inferiore al parietale.
\begin{itemize}
  \item parte interna = \textbf{rocca petrosa} (mediale e in avanti); parte esterna = \textbf{mastoidea} (in basso);
  \item presenta processo zigomatico, fossa glenoidea, processo stiloideo.
\end{itemize}

\subsection{Osso etmoide}
Impari e mediano; parte delle cavità nasali e orbitarie.
\begin{itemize}
  \item articolazioni: frontale (alto), sfenoide (dietro), mascellari e vomere (avanti/basso),
    lacrimali (lati);
  \item due lamine perpendicolari + due masse laterali;
  \item \textbf{lamina cribrosa} (orizzontale): forellini per il nervo olfattivo; faccia superiore
    = parte della fossa cranica, inferiore = volta delle cavità nasali; porta la \textbf{crista
    galli};
  \item masse laterali = labirinto etmoidale; lamina perpendicolare = parte del setto nasale.
\end{itemize}

\subsection{Osso frontale}
Impari e mediano, costituisce la volta cranica.
\begin{itemize}
  \item articolazioni: parietali, etmoide, mascellari, nasali, sfenoide, lacrimali;
  \item parte verticale (\textbf{squama}) + parte orizzontale; margini superiore (parietale),
    inferiore (sopraorbitario), posteriore (sfenoidale);
  \item elementi: sutura metopica, arcata sopracciliare, margine e foro sopraorbitario, cresta
    frontale, incisura etmoidale, seni frontali, fossa orbitaria.
\end{itemize}

\subsection{Osso parietale}
Pari, piatto, quadrangolare irregolare, concavo verso l'interno.
\begin{itemize}
  \item articolazioni: controlaterale (sutura sagittale), frontale (avanti), grande ala dello
    sfenoide e squama del temporale (lati), occipitale (dietro);
  \item faccia esocranica convessa (eminenza parietale); faccia endocranica concava con solchi
    dell'arteria meningea media; linee temporali superiore e inferiore.
\end{itemize}

% ---------------------------------------------------------------------------
\section{Lo splancnocranio}

\subsection{Osso nasale}
Piccolo, laminare, pari, trapezoidale. Articolazioni: frontale (alto), controlaterale (mediale),
processo frontale della mascella (laterale); margine inferiore = contorno superiore dell'apertura
piriforme.

\subsection{Osso lacrimale}
Laminare, pari, quadrangolare. Con la mascella forma la fossa del sacco lacrimale.

\subsection{Osso palatino}
Laminare, pari, tra processo pterigoideo dello sfenoide e superficie posteromediale del mascellare.
Partecipa a cavità nasali, palato e (in parte) pavimento dell'orbita. Porzione verticale +
orizzontale.

\subsection{Cornetto nasale inferiore}
Pari, sotto il cornetto etmoidale medio, nelle cavità nasali; lamina a tegola allungata. Tre
processi superiori:
\begin{itemize}
  \item \textbf{etmoidale}: si unisce al processo uncinato dell'etmoide;
  \item \textbf{lacrimale}: con l'osso lacrimale, chiude il canale nasolacrimale;
  \item \textbf{mascellare}: con la parte inferiore dello iato mascellare.
\end{itemize}

\subsection{Vomere}
Lamina quadrangolare impari, parte del setto nasale. Articolazioni: sfenoide (alto/dietro), lamina
perpendicolare dell'etmoide (alto/avanti), parte orizzontale del palatino e processi palatini delle
mascelle (basso).

\subsection{Osso zigomatico}
Osso breve, pari, quadrangolare.
\begin{itemize}
  \item processo frontale (angolo superiore) $\rightarrow$ osso frontale;
  \item angoli anteriore/inferiore $\rightarrow$ mascella;
  \item processo temporale (angolo posteriore) $\rightarrow$ osso temporale = arco zigomatico;
  \item bordo postero-inferiore $\rightarrow$ grande ala dello sfenoide.
\end{itemize}

\subsection{Osso mascellare}
Pari, il più grande del massiccio facciale; coi controlaterali forma l'arcata mascellare superiore.
\begin{itemize}
  \item corpo + 4 processi: zigomatico, frontale, alveolare, palatino;
  \item corpo piramidale a 4 facce: anteriore (sporgenze sulle radici dei denti, incisura nasale),
    infratemporale/posteriore (separata dall'anteriore dal processo zigomatico).
\end{itemize}

\subsection{Mandibola}
Impari e mediano.
\begin{itemize}
  \item corpo a ferro di cavallo (si articola coi denti dell'arcata inferiore, molari inclusi);
  \item due rami $\rightarrow$ osso temporale; ciascun ramo termina col processo condiloideo
    (condilo) e coronoideo, separati dall'incisura mandibolare;
  \item faccia interna: rilievi/fossette per muscoli estrinseci della lingua e sopraioidei (linea
    miloioidea); foro del canale mandibolare (nervo mandibolare), che emerge dal forame mentoniero
    (faccia esterna, presso la protuberanza mentoniera);
  \item margini: alveolare e basale.
\end{itemize}

% ---------------------------------------------------------------------------
\section{Osso ioide}

\textbf{Osso ioide}: impari e mediano, a ferro di cavallo con concavità posteriore, nella regione
anteriore del collo (in rapporto con cartilagine tiroidea, membrana tiro-ioidea, trachea, muscoli
del pavimento della bocca).
\begin{itemize}
  \item unico osso che \textbf{non si articola} con altre ossa: posizione mantenuta solo dai
    muscoli inseriti;
  \item muscoli: digastrico, stiloioideo, tireoioideo, omoioideo, sternoioideo, costrittore medio
    della faringe, ioglosso, condroglosso, genioglosso, genioioideo.
\end{itemize}

% ---------------------------------------------------------------------------
\section{La colonna vertebrale}

\textbf{Colonna vertebrale} (\textbf{rachide}): principale sostegno del corpo, 33--34 ossa brevi
sovrapposte. Vista di lato, quattro curve alternate:
\begin{itemize}
  \item \textbf{cervicale}: concava posteriormente (lordosi);
  \item \textbf{toracica}: concava anteriormente (cifosi);
  \item \textbf{lombare}: concava posteriormente (lordosi);
  \item \textbf{sacrale}: concava anteriormente (cifosi).
\end{itemize}
Tratti: cervicale, dorsale, lombare, pelvico $\rightarrow$ osso sacro e coccige.

\subsection{Struttura generale delle vertebre}
Ogni \textbf{vertebra} = corpo + arco che delimitano un foro.
\begin{itemize}
  \item \textbf{corpo}: quasi cilindrico, due facce (superiore/inferiore) concave al centro;
  \item vertebre articolate per sinfisi con \textbf{dischi fibrocartilaginei} $\rightarrow$ rendono
    complementari le superfici e ammortizzano i carichi ($\sim$1/3 della lunghezza della colonna);
  \item \textbf{arco vertebrale}: porzione posteriore che completa il forame; due peduncoli, due
    masse apofisarie, due lamine, un processo spinoso.
\end{itemize}

\subsection{Vertebre cervicali (C1--C7)}
Il corpo aumenta di volume in direzione craniocaudale.
\begin{itemize}
  \item \textbf{atlante} (C1): privo di corpo; due archi (anteriore/posteriore) e due masse
    laterali con cavità glenoidea per i condili occipitali; tubercolo anteriore + faccetta per il
    dente; tubercolo posteriore;
  \item \textbf{epistrofeo/asse} (C2): il \textbf{dente} si articola con l'arco anteriore
    dell'atlante (è il corpo dell'atlante rimasto saldato); faccette articolari superiori ai lati;
  \item \textbf{C7} (\textbf{vertebra prominente}): di transizione, processo spinoso lungo, non
    bifido, sporgente;
  \item cervicali tipiche: forame vertebrale, forame trasverso (arteria vertebrale), peduncolo,
    lamina, solco per il nervo spinale, processo uncinato, processo trasverso.
\end{itemize}

\subsection{Vertebre toraciche (T1--T12)}
Si articolano con le coste.
\begin{itemize}
  \item corpi quasi cilindrici (diametri antero-posteriore e trasverso simili), crescenti verso il
    basso;
  \item due faccette articolari per lato (superiore/inferiore) per le teste delle coste; faccetta
    costale trasversaria sul processo trasverso (per il tubercolo della costa);
  \item processi articolari sup./inf., processo spinoso, arco e canale vertebrale.
\end{itemize}

\subsection{Vertebre lombari}
\begin{itemize}
  \item corpo robusto, processo spinoso corto e quadrangolare, ampio forame;
  \item processi trasversi (= costali); processo mammillare sul processo articolare superiore;
  \item da L1 a L5 crescenti; incisure vertebrali sup./inf.\ $\rightarrow$ forame intervertebrale;
    articolazione zigapofisaria tra processi articolari contigui.
\end{itemize}

\subsection{Sacro e coccige}
\begin{itemize}
  \item \textbf{sacro}: fusione di 5 vertebre; impari, piramidale quadrangolare, appiattito;
    si articola con le ossa dell'anca $\rightarrow$ bacino (col coccige);
    \begin{itemize}
      \item faccia anteriore (pelvica) concava e liscia, posteriore (dorsale), due laterali;
      \item 4 coppie di \textbf{forami sacrali} (rami anteriori dei nervi sacrali); linee trasverse
        = residui delle sinostosi;
      \item elementi: promontorio, ala, processo articolare superiore, canale sacrale, creste
        (laterale/mediana/mediale), superficie auricolare per l'anca, tuberosità sacrale, iato
        sacrale, corna, articolazione sacrococcigea;
    \end{itemize}
  \item \textbf{coccige}: fusione di 4--5 segmenti, decrescenti craniocaudalmente.
\end{itemize}

% ---------------------------------------------------------------------------
\section{La gabbia toracica}

\textbf{Coste}: 12 paia (parte ossea + cartilaginea).
\begin{itemize}
  \item \textbf{vere} (1--7): si articolano con lo sterno;
  \item \textbf{asternali} (8, 9, 10);
  \item \textbf{false} (11, 12): non si articolano con lo sterno;
  \item obliquità crescente dalla 1\textsuperscript{a} alla 7\textsuperscript{a}, poi decrescente;
  \item costa = corpo (sottile, piatto, faccia esterna convessa e interna concava, due margini) +
    due estremità; posteriore con testa, collo, tubercolo;
  \item \textbf{1\textsuperscript{a} costa}: raggio di curvatura minore, corta, spessa, appiattita
    (faccia inferiore/superiore); \textbf{2\textsuperscript{a}}: doppia della prima, orientamento
    simile; \textbf{11\textsuperscript{a}--12\textsuperscript{a}}: solo col corpo vertebrale, senza
    collo/tubercolo/angolo/solco costale.
\end{itemize}

\subsection{Sterno}
Osso piatto, impari, mediano; chiude anteriormente il torace ($\sim$20 cm, da T2 a T9). Tre parti:
\begin{itemize}
  \item \textbf{manubrio}: quadrangolare; si articola col corpo tramite sinfisi manubrio-sternale
    (leggeri movimenti in respirazione);
  \item \textbf{corpo} (T3--T8): tra manubrio e processo xifoideo (articolazione xifosternale);
  \item \textbf{processo xifoideo}: la parte più piccola, nell'epigastrio (in rapporto col fegato),
    dà attacco a muscoli.
\end{itemize}

% ---------------------------------------------------------------------------
\section{L'arto superiore}

\textbf{Arto superiore}: spalla, braccio, gomito, avambraccio, mano. Articolazioni a cingolo
scapolare, gomito, polso.
\begin{itemize}
  \item \textbf{cintura pettorale}: scapola + clavicola (spalla);
  \item \textbf{parte libera}: omero, radio e ulna; carpo, metacarpo, falangi.
\end{itemize}

\subsection{Scapola e clavicola}
\begin{itemize}
  \item \textbf{scapola}: osso piatto triangolare pari, posteriore e superolaterale al torace;
    faccia anteriore concava, posteriore divisa dalla \textbf{spina} (continua nell'\textbf{acromion});
    \textbf{cavità glenoidea} (angolo superolaterale) per l'omero;
  \item \textbf{clavicola}: osso piatto a S (convesso avanti nei 2/3 mediali, concavo nel terzo
    laterale); estremità acromiale si articola con l'acromion (artrodia con disco fibrocartilagineo,
    capsula robusta e legamenti).
\end{itemize}

\subsection{Omero}
Osso lungo, scheletro del braccio: diafisi (cilindrica sopra, triangolare sotto) + due epifisi.
\begin{itemize}
  \item epifisi prossimale: testa, collo anatomico, collo chirurgico, tubercolo maggiore e minore
    (separati dal solco intertubercolare), labbri laterale e mediale;
  \item corpo: tuberosità deltoidea; verso il basso si appiattisce $\rightarrow$ \textbf{epicondili}
    laterale e mediale;
  \item \textbf{condilo}: \textbf{capitello} (laterale, per il radio) + \textbf{troclea} (mediale,
    per l'ulna); posteriormente la fossa olecranica.
\end{itemize}
\textbf{Fratture della testa dell'omero}: frequenti negli sport di contatto (sci, snowboard, calcio,
rugby).

\subsection{Radio e ulna}
Entrambi ossa lunghe (due epifisi + diafisi), scheletro dell'avambraccio.
\begin{itemize}
  \item \textbf{ulna} (mediale, sezione triangolare): epifisi prossimale con \textbf{incisura
    trocleare} (per la troclea omerale), processo coronoideo (antero-inferiore) e \textbf{olecrano}
    (postero-superiore); distalmente testa e processo stiloideo;
  \item \textbf{radio} (laterale): diafisi triangolare incurvata, unita all'ulna dalla membrana
    interossea; epifisi prossimale con testa, collo, tuberosità; epifisi distale con processo
    stiloideo (laterale), incisura ulnare (mediale) e superficie articolare concava per il carpo.
\end{itemize}

\subsection{La mano}
\begin{itemize}
  \item \textbf{carpo}: due file $\rightarrow$ prossimale (scafoide/navicolare, semilunare,
    piramidale, pisiforme), distale (trapezio, trapezoide, capitato, uncinato);
  \item \textbf{metacarpo}: 5 ossa lunghe;
  \item \textbf{falangi}: 3 per dito (prossimale, media, distale/ungueale), tranne il primo (2);
    articolate tra loro per troclee, con i metacarpali per condili.
\end{itemize}

% ---------------------------------------------------------------------------
\section{Il bacino}

\textbf{Bacino} (\textbf{pelvi}): anello osseo $\rightarrow$ ossa dell'anca (avanti e lati) + sacro
e coccige (dietro); formazione rigida a cono irregolare, rivestita da tegumenti, muscoli, fasce.
\begin{itemize}
  \item \textbf{cintura pelvica}: due ossa dell'anca unite al sacro dall'\textbf{articolazione
    sacroiliaca};
  \item elementi: promontorio del sacro, cresta iliaca (labbro interno, linea intermedia, labbro
    esterno), ala iliaca, grande e piccola incisura ischiatica, linea arcuata, spina ischiatica,
    eminenza ileo-pubica, rami del pube, foro otturato, tubercolo pubico, arcata pubica.
\end{itemize}

\subsection{Osso dell'anca}
Tre parti:
\begin{itemize}
  \item \textbf{ileo}: corpo (lateralmente l'\textbf{acetabolo}, per il femore) + ala appiattita;
  \item \textbf{ischio}: partecipa all'acetabolo; ramo a ``C'' con la \textbf{tuberosità
    ischiatica};
  \item \textbf{pube}: corpo fuso con ileo e ischio $\rightarrow$ porzione antero-inferiore
    dell'acetabolo.
\end{itemize}
\textbf{Sinfisi pubica}: unione delle porzioni pubiche con disco fibrocartilagineo interposto; ampia,
resistente alle tensioni elastiche.

% ---------------------------------------------------------------------------
\section{L'arto inferiore}

\subsection{Femore}
Osso della coscia; il più lungo, voluminoso e resistente. Corpo (diafisi) + due epifisi.
\begin{itemize}
  \item epifisi prossimale: testa cilindrica, collo anatomico, \textbf{grande} e \textbf{piccolo
    trocantere};
  \item diafisi cilindrica;
  \item epifisi distale: epicondili laterale e mediale; condili fusi avanti (faccia patellare),
    separati dietro dalla fossa intercondilare.
\end{itemize}
\textbf{Fratture}: più frequenti al collo (per il carico), aggravate in età avanzata
dall'\textbf{osteoporosi} (trama ossea rarefatta). La \textbf{vitamina D} è importante per la salute
ossea e la prevenzione dell'osteoporosi.

\subsection{Patella}
Il più grande osso sesamoide; nella regione anteriore del ginocchio, nell'espansione tendinea del
quadricipite femorale, davanti alla faccia patellare del femore. Osso breve triangolare, apice
inferiore; faccia posteriore articolare (porzione mediale e laterale).

\subsection{Tibia e perone}
\begin{itemize}
  \item \textbf{tibia}: osso lungo a sezione triangolare, margine anteriore affilato; \textbf{tuberosità
    tibiale} (tendine del quadricipite); linea del soleo (faccia posteriore); epifisi prossimale a
    piramide con piatto tibiale (per i condili femorali) ed eminenza intercondiloidea; malleolo
    mediale distalmente;
  \item \textbf{perone} (\textbf{fibula}): sottile, laterale alla tibia; testa (prossimale) e
    malleolo laterale (distale); unito alla tibia da membrana interossea e articolazioni
    tibio-fibulari superiore e inferiore.
\end{itemize}

\subsection{Il piede}
\begin{itemize}
  \item \textbf{tarso}: 7 ossa $\rightarrow$ \textbf{talo/astragalo} (testa, collo, corpo, troclea;
    articolazione tibiotarsica), \textbf{calcagno} (sustentaculum tali, tuberosità calcaneale, seno
    del tarso), navicolare, cuboide, tre cuneiformi;
  \item \textbf{metatarso}: 5 ossa lunghe (I--V mediolaterale), incurvate verso il lato plantare;
    prossimalmente artrodie con cuneiformi e cuboide, distalmente condilartrosi con le falangi;
  \item \textbf{falangi}: 3 per dito (2 per il primo), articolate per troclee; le distali accolgono
    l'unghia.
\end{itemize}
